\chapter{La couche limite intégrale : principe et mise en place}\label{ch:cl}
\iterbanner{V0}

\section{Pourquoi une méthode intégrale}
Les équations de Prandtl sont des EDP paraboliques en $(s,y)$. Les résoudre sur chaque face demanderait un maillage dans l'épaisseur, une loi de paroi ou un modèle de turbulence, et des centaines de stations. Les \textbf{méthodes intégrales} font un pari plus modeste : intégrer les équations dans l'épaisseur, ne suivre que quelques grandeurs globales du profil de vitesse, et fermer le système par des corrélations empiriques. On passe d'une EDP à une ou deux \textbf{EDO en $s$}, résolues en quelques millisecondes.

\section{Rappels : grandeurs intégrales et équation de von Kármán}
Pour un profil de vitesse $u(s,y)$ qui tend vers $U_e(s)$ en bord de couche :
\[
\delta^*=\int_0^\infty\Big(1-\frac{u}{U_e}\Big)\dd y,\qquad
\theta=\int_0^\infty\frac{u}{U_e}\Big(1-\frac{u}{U_e}\Big)\dd y,\qquad
H=\frac{\delta^*}{\theta},\qquad c_f=\frac{\tau_w}{\tfrac12\rho U_e^2}.
\]
$\delta^*$ est le déplacement des lignes de courant extérieures, $\theta$ mesure le déficit de quantité de mouvement, et $H$ est un indicateur de la \emph{forme} du profil : environ 2,6 pour un profil laminaire de Blasius, 1,3 à 1,4 pour un profil turbulent plein, et au-delà de 3,5 (laminaire) ou de 2 à 2,8 (turbulent) près du décollement.

L'intégration de l'équation de quantité de mouvement de Prandtl dans l'épaisseur donne l'\textbf{équation intégrale de von Kármán}, valable en laminaire comme en turbulent :
\[
\boxed{\ \frac{\dd\theta}{\dd s}+\big(2+H\big)\frac{\theta}{U_e}\frac{\dd U_e}{\dd s}=\frac{c_f}{2}\ }
\]
Elle se lit comme un bilan : le déficit de quantité de mouvement croît par frottement ($c_f/2$), et un gradient de pression adverse ($U_e'<0$) l'amplifie d'autant plus que le profil est creusé ($H$ grand). Une équation, trois inconnues ($\theta$, $H$, $c_f$) : toutes les méthodes intégrales sont des façons de fermer ce système.

\begin{center}
\begin{tabularx}{\linewidth}{lllX}
\toprule
Régime & Méthode & Inconnues & Fermeture \\
\midrule
laminaire & Thwaites (ch.~\ref{ch:thwaites}) & $\theta$ & $H$ et $c_f$ fonctions de $\lambda=\theta^2U_e'/\nu$ \\
transition & Michel (ch.~\ref{ch:transition}) & --- & seuil sur $\Re_\theta$ \\
turbulent & Head (ch.~\ref{ch:head}) & $\theta$, $H$ & équation d'entraînement, $H_1(H)$, $c_f$ de Ludwieg-Tillmann \\
\bottomrule
\end{tabularx}
\end{center}

\section{De la solution de panneaux aux deux faces}
La couche limite démarre au \textbf{point d'arrêt} et se développe séparément sur l'extrados et sur l'intrados. Il faut donc trouver ce point, puis réorganiser les données de chaque face dans le sens de l'écoulement.

\paragraph{Le point d'arrêt.} En parcourant le contour depuis le bord de fuite extrados, la vitesse tangentielle signée $V_t$ est d'abord négative (l'écoulement remonte le sens de parcours), s'annule au point d'arrêt, puis devient positive sur l'intrados. On cherche l'unique indice $k$ tel que $V_{t,k}<0\le V_{t,k+1}$, et on interpole linéairement le zéro :
\[
s_{\text{arrêt}}=\ell_k+w\,(\ell_{k+1}-\ell_k),\qquad w=\frac{V_{t,k}}{V_{t,k}-V_{t,k+1}},
\]
où $\ell_k$ est l'abscisse curviligne cumulée du point de collocation $k$ le long du contour.

\paragraph{Les deux faces.} Chaque face reçoit ensuite trois tableaux, orientés du point d'arrêt vers le bord de fuite :
\begin{itemize}
\item l'abscisse curviligne $s$ mesurée depuis le point d'arrêt (donc $s>0$) ;
\item l'abscisse $x/c$ des stations, pour repérer la transition et la coupure en $x_{stop}$ ;
\item la vitesse extérieure $U_e=|V_t|$ : $-V_t$ sur l'extrados, $+V_t$ sur l'intrados.
\end{itemize}

\begin{figure}[htbp]
\centering
% figure générée par tools/make_figures.py : vt_contour
\begin{tikzpicture}
\begin{groupplot}[group style={group size=2 by 1,horizontal sep=1.6cm},width=0.5\linewidth,height=5.4cm,
  title style={font=\small},tick label style={font=\scriptsize},legend style={draw=none,font=\scriptsize}]
\nextgroupplot[title={$V_t$ signée le long du contour},xlabel={abscisse $\ell$ depuis le BF extrados},ylabel={$V_t/U_\infty$},xmin=0,xmax=2.042]
\addplot[cink,very thick] coordinates {(0.00012493,-0.69892) (0.00062451,-0.75574) (0.0016232,-0.79253) (0.0031199,-0.81972) (0.0051131,-0.84166) (0.0076009,-0.8603) (0.010581,-0.8767) (0.01405,-0.89146) (0.018004,-0.90497) (0.02244,-0.91751) (0.027352,-0.92927) (0.032737,-0.94037) (0.038589,-0.95092) (0.044901,-0.96101) (0.051668,-0.97069) (0.058882,-0.98001) (0.066537,-0.98901) (0.074624,-0.99773) (0.083136,-1.0062) (0.092064,-1.0144) (0.1014,-1.0224) (0.11113,-1.0303) (0.12125,-1.0379) (0.13175,-1.0454) (0.14262,-1.0527) (0.15385,-1.0599) (0.16542,-1.0669) (0.17732,-1.0739) (0.18955,-1.0807) (0.20208,-1.0874) (0.21492,-1.0941) (0.22804,-1.1006) (0.24143,-1.1071) (0.25508,-1.1136) (0.26898,-1.1199) (0.28311,-1.1263) (0.29746,-1.1325) (0.312,-1.1388) (0.32674,-1.145) (0.34166,-1.1512) (0.35673,-1.1574) (0.37195,-1.1636) (0.3873,-1.1698) (0.40276,-1.176) (0.41832,-1.1822) (0.43396,-1.1884) (0.44968,-1.1946) (0.46544,-1.2009) (0.48125,-1.2072) (0.49707,-1.2136) (0.5129,-1.22) (0.52872,-1.2265) (0.54452,-1.2331) (0.56027,-1.2398) (0.57597,-1.2468) (0.59159,-1.2542) (0.60717,-1.2621) (0.6227,-1.27) (0.63814,-1.277) (0.65345,-1.2835) (0.66862,-1.2895) (0.68362,-1.2952) (0.69844,-1.3007) (0.71308,-1.3058) (0.72751,-1.3107) (0.74172,-1.3153) (0.7557,-1.3197) (0.76943,-1.3238) (0.78291,-1.3277) (0.79611,-1.3313) (0.80903,-1.3346) (0.82166,-1.3377) (0.83398,-1.3406) (0.84598,-1.3432) (0.85766,-1.3456) (0.86899,-1.3477) (0.87998,-1.3496) (0.89061,-1.3512) (0.90087,-1.3525) (0.91076,-1.3535) (0.92026,-1.3543) (0.92938,-1.3546) (0.9381,-1.3546) (0.94642,-1.3541) (0.95433,-1.3531) (0.96183,-1.3512) (0.96892,-1.3484) (0.97559,-1.3442) (0.98185,-1.3383) (0.9877,-1.3297) (0.99314,-1.3175) (0.99818,-1.3001) (1.0028,-1.2752) (1.0071,-1.2393) (1.011,-1.1877) (1.0146,-1.1141) (1.0179,-1.0107) (1.021,-0.87072) (1.0239,-0.69135) (1.0267,-0.47892) (1.0295,-0.25021) (1.0323,-0.022469) (1.0352,0.18632) (1.0383,0.36477) (1.0416,0.50981) (1.0452,0.62418) (1.0491,0.7131) (1.0533,0.78198) (1.058,0.8355) (1.0629,0.87732) (1.0683,0.91023) (1.0741,0.93632) (1.0803,0.95713) (1.0869,0.97381) (1.0938,0.98722) (1.1012,0.99802) (1.109,1.0067) (1.1172,1.0136) (1.1257,1.0191) (1.1346,1.0235) (1.1439,1.0268) (1.1536,1.0292) (1.1636,1.031) (1.174,1.0321) (1.1847,1.0326) (1.1957,1.0328) (1.2071,1.0325) (1.2188,1.0318) (1.2308,1.0309) (1.243,1.0297) (1.2556,1.0283) (1.2684,1.0267) (1.2815,1.025) (1.2948,1.0232) (1.3084,1.0212) (1.3222,1.0193) (1.3362,1.0173) (1.3504,1.0153) (1.3648,1.0133) (1.3793,1.0114) (1.3941,1.0097) (1.4089,1.0081) (1.4239,1.0068) (1.4391,1.0061) (1.4544,1.0056) (1.4698,1.0047) (1.4853,1.0034) (1.5009,1.0018) (1.5165,1.0001) (1.5321,0.99823) (1.5477,0.99633) (1.5634,0.99438) (1.579,0.99241) (1.5946,0.99041) (1.6102,0.9884) (1.6256,0.98639) (1.641,0.98439) (1.6564,0.98238) (1.6716,0.98039) (1.6866,0.9784) (1.7016,0.97641) (1.7164,0.97443) (1.731,0.97246) (1.7454,0.97048) (1.7596,0.9685) (1.7737,0.9665) (1.7875,0.9645) (1.801,0.96246) (1.8143,0.9604) (1.8274,0.9583) (1.8401,0.95616) (1.8526,0.95396) (1.8648,0.95169) (1.8766,0.94935) (1.8881,0.94692) (1.8993,0.9444) (1.9101,0.94176) (1.9206,0.93899) (1.9307,0.93608) (1.9404,0.93301) (1.9497,0.92976) (1.9586,0.92631) (1.9671,0.92264) (1.9752,0.91871) (1.9828,0.9145) (1.99,0.90997) (1.9968,0.90507) (2.0031,0.89974) (2.0089,0.89392) (2.0143,0.88753) (2.0192,0.88044) (2.0236,0.87252) (2.0276,0.86357) (2.0311,0.8533) (2.034,0.8413) (2.0365,0.82692) (2.0385,0.80909) (2.04,0.78577) (2.041,0.75256) (2.0415,0.69892)};
\addplot[cmuted,thin] coordinates {(0,0) (2.042,0)};
\addplot[only marks,mark=*,corange,mark size=2.4pt] coordinates {(1.0326,0)};
\node[corange,font=\scriptsize,anchor=north west] at (axis cs:1.063,-0.08) {point d'arrêt};
\node[font=\scriptsize,anchor=north] at (axis cs:0.5,-0.2) {extrados ($V_t<0$)};
\node[font=\scriptsize,anchor=north] at (axis cs:1.6,0.8) {intrados ($V_t>0$)};
\nextgroupplot[title={$U_e=|V_t|$ sur chaque face},xlabel={$x/c$},ylabel={$U_e/U_\infty$},xmin=0,xmax=1,legend pos=south east]
\addplot[cblue,very thick] coordinates {(0.0010289,0.022469) (0.00026156,0.25021) (-1.4838e-05,0.47892) (0.00020448,0.69135) (0.00092198,0.87072) (0.0021378,1.0107) (0.0038514,1.1141) (0.0060621,1.1877) (0.0087685,1.2393) (0.011969,1.2752) (0.01566,1.3001) (0.01984,1.3175) (0.024505,1.3297) (0.029651,1.3383) (0.035273,1.3442) (0.041368,1.3484) (0.047928,1.3512) (0.054949,1.3531) (0.062423,1.3541) (0.070344,1.3546) (0.078703,1.3546) (0.087494,1.3543) (0.096707,1.3535) (0.10633,1.3525) (0.11636,1.3512) (0.12679,1.3496) (0.13759,1.3477) (0.14877,1.3456) (0.16031,1.3432) (0.17219,1.3406) (0.18442,1.3377) (0.19696,1.3346) (0.20982,1.3313) (0.22297,1.3277) (0.2364,1.3238) (0.25011,1.3197) (0.26407,1.3153) (0.27826,1.3107) (0.29268,1.3058) (0.30731,1.3007) (0.32214,1.2952) (0.33714,1.2895) (0.3523,1.2835) (0.36761,1.277) (0.38305,1.27) (0.39857,1.2621) (0.41414,1.2542) (0.42974,1.2468) (0.44542,1.2398) (0.46115,1.2331) (0.47691,1.2265) (0.49269,1.22) (0.50848,1.2136) (0.52425,1.2072) (0.54,1.2009) (0.5557,1.1946) (0.57135,1.1884) (0.58692,1.1822) (0.6024,1.176) (0.61777,1.1698) (0.63302,1.1636) (0.64814,1.1574) (0.6631,1.1512) (0.6779,1.145) (0.69252,1.1388) (0.70695,1.1325) (0.72116,1.1263) (0.73515,1.1199) (0.74891,1.1136) (0.76242,1.1071) (0.77566,1.1006) (0.78863,1.0941) (0.80131,1.0874) (0.81369,1.0807) (0.82576,1.0739) (0.8375,1.0669) (0.84891,1.0599) (0.85997,1.0527) (0.87068,1.0454) (0.88102,1.0379) (0.89098,1.0303) (0.90055,1.0224) (0.90973,1.0144) (0.9185,1.0062) (0.92686,0.99773) (0.9348,0.98901) (0.94232,0.98001) (0.9494,0.97069) (0.95603,0.96101) (0.96222,0.95092) (0.96795,0.94037) (0.97323,0.92927) (0.97804,0.91751) (0.98238,0.90497) (0.98625,0.89146) (0.98965,0.8767) (0.99257,0.8603) (0.995,0.84166) (0.99695,0.81972) (0.99841,0.79253) (0.99939,0.75574) (0.99988,0.69892)};\addlegendentry{extrados}
\addplot[corange,very thick] coordinates {(0.0022837,0.18632) (0.0040239,0.36477) (0.006247,0.50981) (0.0089501,0.62418) (0.01213,0.7131) (0.015782,0.78198) (0.019902,0.8355) (0.024485,0.87732) (0.029527,0.91023) (0.035021,0.93632) (0.040961,0.95713) (0.047342,0.97381) (0.054155,0.98722) (0.061395,0.99802) (0.069053,1.0067) (0.077121,1.0136) (0.085592,1.0191) (0.094457,1.0235) (0.10371,1.0268) (0.11333,1.0292) (0.12333,1.031) (0.13367,1.0321) (0.14437,1.0326) (0.1554,1.0328) (0.16676,1.0325) (0.17844,1.0318) (0.19042,1.0309) (0.20269,1.0297) (0.21525,1.0283) (0.22808,1.0267) (0.24116,1.025) (0.25449,1.0232) (0.26806,1.0212) (0.28185,1.0193) (0.29585,1.0173) (0.31005,1.0153) (0.32443,1.0133) (0.33899,1.0114) (0.3537,1.0097) (0.36855,1.0081) (0.38354,1.0068) (0.39866,1.0061) (0.41396,1.0056) (0.42937,1.0047) (0.44486,1.0034) (0.4604,1.0018) (0.47599,1.0001) (0.4916,0.99823) (0.50723,0.99633) (0.52285,0.99438) (0.53845,0.99241) (0.55402,0.99041) (0.56954,0.9884) (0.58499,0.98639) (0.60037,0.98439) (0.61565,0.98238) (0.63082,0.98039) (0.64587,0.9784) (0.66078,0.97641) (0.67553,0.97443) (0.69012,0.97246) (0.70452,0.97048) (0.71872,0.9685) (0.73272,0.9665) (0.74649,0.9645) (0.76001,0.96246) (0.77329,0.9604) (0.7863,0.9583) (0.79903,0.95616) (0.81147,0.95396) (0.82361,0.95169) (0.83542,0.94935) (0.84691,0.94692) (0.85806,0.9444) (0.86886,0.94176) (0.8793,0.93899) (0.88936,0.93608) (0.89904,0.93301) (0.90832,0.92976) (0.9172,0.92631) (0.92567,0.92264) (0.93372,0.91871) (0.94134,0.9145) (0.94852,0.90997) (0.95526,0.90507) (0.96155,0.89974) (0.96737,0.89392) (0.97274,0.88753) (0.97763,0.88044) (0.98205,0.87252) (0.986,0.86357) (0.98945,0.8533) (0.99242,0.8413) (0.9949,0.82692) (0.99689,0.80909) (0.99838,0.78577) (0.99938,0.75256) (0.99988,0.69892)};\addlegendentry{intrados}
\end{groupplot}
\end{tikzpicture}

\caption{NACA 2412, $\alpha=2^\circ$. À gauche, la vitesse tangentielle signée le long du contour, parcouru depuis le bord de fuite extrados. Elle s'annule au point d'arrêt, situé ici en $x/c\approx0{,}001$ sur l'intrados. À droite, les vitesses extérieures que reçoit chaque face.}\label{fig:vt}
\end{figure}

\begin{codebox}[ : \code{boundary\_layer.split\_at\_stagnation}]
\begin{lstlisting}
k = np.nonzero((Vt[:-1] < 0) & (Vt[1:] >= 0))[0]    # changement de signe - -> +
if k.size != 1: raise RuntimeError("point d'arret ambigu")
s_nodes = np.concatenate([[0.0], np.cumsum(P.L)])
s_c = 0.5 * (s_nodes[:-1] + s_nodes[1:])          # abscisse des collocations
w = Vt[k] / (Vt[k] - Vt[k + 1])
s_stag = s_c[k] + w * (s_c[k + 1] - s_c[k])
up = np.arange(k, -1, -1)       # du point d'arret vers le BF extrados
lo = np.arange(k + 1, P.n)      # du point d'arret vers le BF intrados
upper = (s_stag - s_c[up], P.xc[up], -Vt[up], up)
lower = (s_c[lo] - s_stag, P.xc[lo],  Vt[lo], lo)
\end{lstlisting}
Si le signe change plusieurs fois (forme pathologique, très forte incidence), l'exception est rattrapée par \code{evaluate}, qui renvoie un résultat marqué invalide au lieu de planter.
\end{codebox}

\section{L'organisation de la marche sur une face}
La fonction \code{march\_surface(s, x, Ue, Re, ...)} traite une face en cinq temps :
\begin{enumerate}
\item couper la face en $x_{stop}=0{,}985$ (ch.~\ref{ch:trainee} explique pourquoi) ;
\item calculer la solution laminaire de Thwaites sur \emph{toute} la face, ce qui ne coûte qu'une somme cumulée ;
\item chercher la première station où un critère de transition est franchi, et interpoler sa position ;
\item intégrer les équations de Head de la transition jusqu'au bout ;
\item renvoyer un objet \code{SurfaceBL} avec $s$, $x$, $U_e$, $\theta$, $H$, $c_f$, un masque turbulent, $x_{tr}$, le mode de transition, $x_{sep}$ et des drapeaux.
\end{enumerate}
Les valeurs au bout de la face, $\theta_{TE}$, $H_{TE}$ et $U_{e,TE}$, alimentent Squire-Young.

\begin{piege}
$U_e$ est mis à $\max(U_e,10^{-6})$ avant tout calcul. La première station est très proche du point d'arrêt, où $U_e\to0$, et plusieurs formules divisent par $U_e$ ou par $U_e^6$. Le seuil évite une division par zéro sans changer le résultat.
\end{piege}

\exercices
\begin{exercice}
Calculer $\delta^*/\delta$, $\theta/\delta$ et $H$ pour le profil linéaire $u/U_e=y/\delta$ ($0\le y\le\delta$), puis pour le profil en puissance $u/U_e=(y/\delta)^{1/7}$. Lequel ressemble à un profil laminaire, lequel à un profil turbulent ?
\end{exercice}
\begin{exercice}
Pourquoi le point d'arrêt d'un profil cambré à incidence positive est-il sur l'intrados, légèrement en aval du bord d'attaque géométrique ? Dans quel sens se déplace-t-il quand $\alpha$ augmente ?
\end{exercice}
